\documentclass{beamer}
\usetheme{Warsaw}
\graphicspath{ {./images/} }

\title[Binding Residue Prediction]{Protein-Ligand Binding Residue Prediction using CNN}
\author{
                Abdus Salam Khazi\\
                \href{mailto:abdus.khazi@students.uni-freiburg.de}
                        {abdus.khazi@students.uni-freiburg.de}
            }
\institute{Albert-Ludwigs-Universität Freiburg}
\date{\today}
\logo{\includegraphics[width=.2\textwidth]{Logo}}

\hypersetup{
    colorlinks=true,
    linkcolor=blue,
    filecolor=magenta,      
    urlcolor=cyan,
}

\begin{document}

\begin{frame}
\titlepage
\end{frame}

\begin{frame}[t]{Biological Background}
What are protein? \pause
\begin{itemize}
\item They are complex molecules that are work horses of a living organism.
They are literally machines.
\href{https://www.youtube.com/watch?v=zm-3kovWpNQ}{TEDx Talk}
\pause

\item Proteins are essential for almost every function of a living organism body.
E.g: Movement of muscles, Transportation of oxygen, Digestion, Immune system etc.
\pause

\item Below are 2 proteins - Hemoglobin(Left) and Motor Protiens(Right).

\includegraphics[scale=0.25]{Hemoglobin}
\includegraphics[scale=0.25]{MotorProtein}
\end{itemize}
\end{frame}



\begin{frame}[t]{Biological Background}
Protein formation process.
\href{https://www.youtube.com/watch?v=gG7uCskUOrA}{DNA to Protein}
\pause
\begin{itemize}
\item All Proteins required by an organism are encodied in the DNA.
\pause
\item First, RNA polymerase transcribes the part of our DNA which contains protein encoding to messenger RNA a.k.a mRNA.
\pause
\item  mRNA is post processed to get an RNA. This RNA goes outside the nucleas of a cell. RNA $\in \{A, T, C, G\}^n$ where $n \in \mathbb{I}^+$
\pause
\item  Ribosomes in the cell cytoplasm then translate RNA to an amino acide sequence. This sequence is called primary structure of a protein. The mapping can be written as - $Ribosome:\mathbb \{A, T, C, G\}^n \to \{R, N,D, C ...\}^m$
where $ m \approx \frac{n}{3}$. 
\pause
\item This amino acid sequence folds into a 3D shape depending on its sequence. This structure is called tertiary protien structure.
\end{itemize}
\end{frame}

\begin{frame}[t]{Biological Background}
What are Ligands?
\pause
\begin{itemize}
\item They are molecules that bind to particular proteins, called receptor proteins.
They act as signals to a binding protein, which may intern start a cascading effect.
\item Example - Glucagon is a ligand created by pancreas. This binds to Glucagon receptor protein which signals it to control the blood sugar level.
\end{itemize}

\pause
Why study protein Ligand Interaction?
\pause

\begin{itemize}
\item It will help us understand biological processes better
\item In the above example we can develop better treatment for type-2 diabetes by understanding the working of Glucaogon and its receptor protein.
\end{itemize}

\end{frame}

\end{document}